\documentclass[border=0pt]{standalone}
\newcommand{\revlamina}{1}
\providecommand{\revlamina}{1}
% Preámbulo común de las láminas VE (A3 horizontal, 420 x 297 mm)
\usepackage[T1]{fontenc}
\usepackage[utf8]{inputenc}
\usepackage[spanish,es-nodecimaldot]{babel}
\usepackage[scaled=0.92]{helvet}
\renewcommand{\familydefault}{\sfdefault}
\usepackage{tikz}
\usetikzlibrary{arrows.meta,patterns,calc,babel,positioning,decorations.pathreplacing}
\usepackage{siunitx}
\sisetup{output-decimal-marker={,},mode=text}
\DeclareSIUnit\kWp{kWp}
\DeclareSIUnit\ohm{\ensuremath{\Omega}}

\definecolor{inv1}{RGB}{37,99,180}
\definecolor{inv2}{RGB}{22,128,84}
\definecolor{nfpa}{RGB}{200,30,30}
\definecolor{viento}{RGB}{40,70,200}
\definecolor{sombra}{RGB}{230,130,0}
\definecolor{losa}{RGB}{120,190,200}
\definecolor{dc}{RGB}{225,95,0}
\definecolor{exist}{RGB}{110,110,110}
\definecolor{nuevo}{RGB}{190,20,20}

% Marcador de referencia numerada
\newcommand{\refm}[2]{\node[circle,draw=black,fill=white,line width=0.35pt,inner sep=0pt,minimum size=3.4mm,font=\fontsize{5.5}{6}\selectfont\bfseries] at #1 {#2};}

% Cajetín: \cajetin{lamina}{titulo}{contenido}{escala}{hoja}
\newcommand{\cajetin}[5]{%
  \draw[line width=0.6pt] (250,10) rectangle (410,58);
  \draw[line width=0.3pt] (250,46) -- (410,46);
  \draw[line width=0.3pt] (250,34) -- (410,34);
  \draw[line width=0.3pt] (250,22) -- (410,22);
  \draw[line width=0.3pt] (370,10) -- (370,34);
  \draw[line width=0.3pt] (330,10) -- (330,22);
  \node[anchor=north west,font=\fontsize{6}{7}\selectfont,align=left] at (251.5,57) {\textbf{Instituto Tecnológico de Costa Rica} -- EL-4601 Normalización Técnica para Electrónica\\Proyecto No.~3: Cargadores para vehículos eléctricos -- diseño en Costa Rica};
  \node[anchor=north west,font=\fontsize{6}{7}\selectfont,align=left] at (251.5,45) {\textbf{Proyecto:} Estación de carga VE de uso público -- Parqueo norte, campus UTN Villa Bonita (CAA-TEC)\\\textbf{Propietario:} Universidad Técnica Nacional -- Catastro A-652521-2000};
  \node[anchor=north west,font=\fontsize{6.5}{7.5}\selectfont,align=left,text width=115mm] at (251.5,33) {\textbf{Contenido:} #2\\#3};
  \node[anchor=north west,font=\fontsize{5.5}{6.5}\selectfont,align=left] at (371.5,33) {\textbf{Lámina}};
  \node[font=\fontsize{16}{18}\selectfont\bfseries] at (390,25.5) {#1};
  \node[anchor=north west,font=\fontsize{5.5}{6.5}\selectfont,align=left] at (251.5,21) {\textbf{Elaboró:} Grupo de trabajo EL-4601\\\textbf{Revisó:} \rule{25mm}{0.2pt}\\\textbf{Fecha:} 3/11/2026 \quad \textbf{Rev.:} \revlamina};
  \node[anchor=north west,font=\fontsize{5.5}{6.5}\selectfont,align=left] at (331.5,21) {\textbf{Escala (A3):}\\#4};
  \node[anchor=north west,font=\fontsize{5.5}{6.5}\selectfont,align=left] at (371.5,21) {\textbf{Hoja:}\\#5};
}
% Marco de la lámina
\newcommand{\marco}{%
  \draw[line width=0.2pt,white] (0,0) rectangle (420,297);
  \draw[line width=0.8pt] (10,10) rectangle (410,287);
}

\definecolor{ac}{RGB}{37,99,180}
\definecolor{gbt}{RGB}{22,128,84}
\definecolor{dcc}{RGB}{225,95,0}
\definecolor{ali}{RGB}{190,20,20}
\definecolor{tierra}{RGB}{30,140,60}

\tikzset{
  txt/.style={font=\fontsize{5.4}{6.5}\selectfont,align=center},
  tl/.style={font=\fontsize{5.4}{6.5}\selectfont,align=left},
  cota/.style={<->,>={Latex[length=1.3mm]},line width=0.25pt},
  conc/.style={fill=gray!30,draw=black,line width=0.4pt},
  tit/.style={font=\fontsize{7.5}{9}\selectfont\bfseries,anchor=north west},
}
\newcommand{\cv}[4]{\draw[cota] (#3,#1) -- node[fill=white,inner sep=0.4pt,font=\fontsize{4.8}{5.5}\selectfont,rotate=90]{#4} (#3,#2);}
\newcommand{\ch}[4]{\draw[cota] (#1,#3) -- node[fill=white,inner sep=0.4pt,font=\fontsize{4.8}{5.5}\selectfont]{#4} (#2,#3);}
\begin{document}
\begin{tikzpicture}[x=1mm,y=1mm]
\marco
\node[anchor=north west,font=\fontsize{10}{12}\selectfont\bfseries] at (15,283) {DETALLES DE MONTAJE, CANALIZACIÓN, PUESTA A TIERRA, ROTULACIÓN Y EMERGENCIA};

% =============== D1: pedestal CA (elevación, 1:20 -> 1 m = 50 mm)
\node[tit] at (15,274) {D1 -- PEDESTAL DOBLE DE CARGA CA (PED-1/PED-2), ELEVACIÓN 1:20};
\begin{scope}[shift={(25,175)},x=50mm,y=50mm]
\fill[gray!15] (-0.2,-0.3) rectangle (1.6,0);          % terreno
\draw[line width=0.5pt] (-0.2,0) -- (1.6,0);
\draw[conc] (0.45,0) rectangle (1.05,0.15);              % isla 150 mm
\draw[conc,pattern=crosshatch dots,pattern color=gray] (0.55,-0.25) rectangle (0.95,0);   % dado
\draw[line width=0.5pt,fill=gray!50] (0.71,0.15) rectangle (0.79,1.55);   % poste
\draw[line width=0.5pt,fill=ac!60] (0.79,0.95) rectangle (0.86,1.25);    % cargador
\draw[line width=0.5pt,fill=ac!60] (0.64,0.95) rectangle (0.71,1.25);
\draw[line width=0.4pt,fill=black] (0.86,0.98) rectangle (0.9,1.06);      % funda conector
\draw[ac,line width=0.8pt] (0.86,1.0) .. controls (1.15,0.9) and (1.2,0.5) .. (1.05,0.35) .. controls (0.95,0.25) and (0.9,0.6) .. (0.9,1.02);
\draw[line width=0.5pt] (0.75,1.55) -- (0.75,1.62) -- (1.05,1.62);       % brazo gestor de cable
\draw[gray,line width=0.3pt] (1.05,1.62) -- (1.05,1.2);
\node[tl,anchor=west] at (1.08,1.4) {gestor de cable\\(625.17)};
\draw[line width=0.4pt,fill=yellow!80!black] (0.05,0.15) rectangle (0.13,1.0);   % bolardo
\draw[conc] (0.0,-0.3) rectangle (0.18,0.0);
\draw[ac!70!black,line width=0.8pt] (0.75,-0.3) -- (0.75,0.15);
\node[tl,anchor=west] at (0.78,-0.18) {PVC 1'' (2 circuitos)};
\fill[gray!60] (1.25,0) rectangle (1.55,0.12);  % tope rueda
\cv{0}{1.0}{-0.1}{1,00}
\cv{0.15}{1.02}{0.4}{0,87}
\cv{0}{1.25}{1.5}{1,25 máx}
\ch{0.13}{0.71}{1.7}{$\ge$ 0,60}
\node[txt,anchor=north] at (0.09,-0.32) {bolardo \O{}150\\relleno};
\node[txt,anchor=north] at (0.75,-0.32) {dado 400$\times$400\\$\times$250, anclas\\4 $\times$ 13 mm};
\node[txt,anchor=north] at (1.4,-0.02) {tope};
\node[tl,anchor=west] at (0.92,1.1) {AC-1/AC-2\\NEMA 4};
\end{scope}
\node[tl,anchor=north west,text width=88mm] at (15,151) {Conector alojado entre 0,9 y 1,2 m sobre el piso terminado (alcance de silla de ruedas, Ley 7600). Cable de 7,5 m recogido en el gestor de cable sin tocar el piso (NEC 625.17). Anclaje al dado con pernos de acero inoxidable (corrosión, NEC 300.6). Montaje según el manual del fabricante (110.3(B)).};

% =============== D2: zanja (1:10 -> 1 m = 100 mm)
\node[tit] at (118,274) {D2 -- SECCIÓN DE ZANJA DEL ALIMENTADOR AL-VE, 1:15};
\begin{scope}[shift={(128,262)},x=66.67mm,y=66.67mm]
\fill[brown!20] (-0.1,-0.7) rectangle (0.6,0);
\fill[white] (0,-0.7) rectangle (0.45,0);
\fill[yellow!12] (0,-0.7) rectangle (0.45,-0.42);          % arena
\fill[brown!10] (0,-0.42) rectangle (0.45,-0.05);          % relleno compactado
\fill[gray!70] (-0.1,-0.05) rectangle (0.6,0);             % asfalto
\draw[line width=0.5pt] (0,0) -- (0,-0.7) -- (0.45,-0.7) -- (0.45,0);
\foreach \x in {0.13,0.32}{\draw[line width=0.6pt,fill=white] (\x,-0.6) circle (0.0445);}
\fill[ali] (0.13,-0.6) circle (0.012); \fill[ali] (0.12,-0.585) circle (0.012); \fill[ali] (0.14,-0.585) circle (0.012); \fill[tierra] (0.13,-0.62) circle (0.007);
\draw[red,line width=0.8pt,dashed] (0.02,-0.3) -- (0.43,-0.3);
\cv{-0.6}{0}{-0.05}{0,60}
\cv{-0.3}{0}{0.52}{0,30}
\ch{0}{0.45}{-0.76}{0,45}
\node[tl,anchor=west] at (0.6,-0.02) {asfalto repuesto 50 mm};
\node[tl,anchor=west] at (0.6,-0.2) {relleno compactado al 95 \%};
\node[tl,anchor=west] at (0.6,-0.3) {\textcolor{red}{cinta de advertencia ``PELIGRO -- 480 V''}};
\node[tl,anchor=west] at (0.6,-0.52) {cama y recubrimiento de arena 100 mm};
\node[tl,anchor=west] at (0.6,-0.62) {AL-VE: PVC 3'' cédula 40 (3\#3/0 + N\#4 + EGC\#6)\\y ducto PVC 3'' de reserva};
\end{scope}
\node[tl,anchor=north west,text width=105mm] at (118,206) {Profundidad mínima 600 mm bajo calles y parqueos para tubería no metálica (NEC 300.5, Tabla 300.5). Registros de 60$\times$60 cm cada $\le$ 30 m y en cada cambio de dirección; sellado de ductos contra el ingreso de agua (300.5(G)).};

% =============== D3: base de DC-1 (planta 1:25 -> 1 m = 40 mm)
\node[tit] at (237,274) {D3 -- BASE DEL CARGADOR RÁPIDO DC-1, PLANTA 1:25};
\begin{scope}[shift={(258,205)},x=40mm,y=40mm]
\draw[conc] (0,0) rectangle (1.4,1.3);
\draw[line width=0.6pt,fill=dcc!60] (0.29,0.45) rectangle (1.11,1.05) node[midway,txt]{DC-1\\820 $\times$ 600};
\foreach \x/\y in {0.36/0.52,1.04/0.52,0.36/0.98,1.04/0.98}{\draw (\x,\y) circle (0.02);}
\draw[dcc!80!black,line width=0.6pt] (0.7,0.75) circle (0.05);
\node[tl,anchor=west] at (1.45,0.75) {entrada inferior\\PVC 1-1/2''};
\foreach \x in {0.1,1.3}{\node[circle,draw=black,fill=yellow!80!black,inner sep=0pt,minimum size=2.5mm] at (\x,-0.25) {};}
\draw[gray,dashed] (0,-0.25) -- (1.4,-0.25);
\ch{0}{1.4}{-0.45}{1,40}
\cv{0}{1.3}{-0.15}{1,30}
\cv{-0.25}{0.45}{1.55}{$\ge$ 0,70}
\node[tl,anchor=north west] at (0,-0.55) {frente (conectores) hacia P5/P6 $\rightarrow$};
\end{scope}
\node[tl,anchor=north west,text width=150mm] at (237,163) {Losa de concreto de 200 mm, 150 mm sobre la isla, con 4 anclas de acero inoxidable según la plantilla del fabricante. Gabinete NEMA 3R / IK10, 1950 mm de alto; dejar $\ge$ 0,7 m libres al frente y los espacios de ventilación del manual. Cables CCS1 y CHAdeMO de 4,6 m con soporte; botón de emergencia integrado.};

% =============== D4: puesta a tierra
\node[tit] at (15,133) {D4 -- PUESTA A TIERRA DE LA CELDA ELÉCTRICA (sin escala)};
\begin{scope}[shift={(20,85)}]
\fill[brown!12] (0,-30) rectangle (95,0);
\draw[line width=0.4pt] (0,0) -- (95,0);
\draw[conc] (12,0) rectangle (72,4);
\node[txt] at (42,2) {losa de la celda};
\foreach \x in {15,70}{\draw[tierra,line width=1.4pt] (\x,-2) -- (\x,-28); \draw[line width=0.4pt,fill=white] (\x-2.5,-4) rectangle (\x+2.5,0);}
\draw[tierra,line width=0.8pt] (15,-2) -- (70,-2);
\node[txt,anchor=north] at (42.5,-2.6) {\#4 Cu desnudo, soldadura exotérmica};
\draw[tierra,line width=0.8pt] (30,-2) -- (30,12) node[txt,anchor=south]{GEC \#6\\TD-VE};
\draw[tierra,line width=0.8pt] (52,-2) -- (52,12) node[txt,anchor=south]{GEC \#6\\T-VE (SDS)};
\node[txt,anchor=west] at (73,-15) {2 varillas Cu\\16 mm $\times$ 2,4 m\\separación $\ge$ 1,8 m};
\ch{15}{70}{-31.5}{$\ge$ 1,8 m (250.53(A)(3))}
\node[txt,anchor=north] at (15,-28.5) {};
\end{scope}
\node[tl,anchor=north west,text width=100mm] at (15,49) {Electrodo de la estructura alimentada por AL-VE (250.32(A)); dos varillas (250.53(A)(2)). Registros de inspección con tapa. Medir la resistencia en la puesta en marcha; objetivo de diseño $\le$ 10 $\Omega$ (S). Pedestales, bolardos metálicos, luminarias y la malla del cerramiento se unen al EGC.};

% =============== D5: rotulación
\node[tit] at (122,133) {D5 -- ROTULACIÓN};
\node[anchor=north west,font=\fontsize{5.2}{6.3}\selectfont] at (122,126) {%
\begin{tabular}{@{}p{19mm}p{84mm}@{}}
\textbf{Ubicación} & \textbf{Texto / norma}\\\hline
TD-VE & ``TABLERO ESTACIÓN VE -- 480Y/277 V -- ALIMENTADO DESDE IP-VE, NICHO TRANSF. ICE 119688''; corriente de falla disponible 5,8 kA y fecha (110.24); directorio de circuitos (408.4)\\
IP-VE (nicho) & ``ESTACIÓN DE CARGA VE -- DESCONECTADOR PRINCIPAL'' (230.70(B)); riesgo de arco eléctrico (110.16)\\
Interruptor DC-1 & ``MEDIO DE DESCONEXIÓN DC-1 -- BLOQUEABLE'' (625.43)\\
T-VE & ``SISTEMA DERIVADO 208Y/120 V'' (250.30); ``PELIGRO -- 480 V''\\
EPO-1/EPO-2 & ``PARO DE EMERGENCIA -- DESENERGIZA TODOS LOS CARGADORES''\\
Cada EVSE & Identificador de compatibilidad INTE N120 (hexágono CA / cuadrado CC); potencia; teléfono de soporte; instrucciones de uso\\
P1 & Símbolo internacional de accesibilidad, señal vertical y pintura (Ley 7600)\\
P1--P6 & ``SOLO VEHÍCULOS ELÉCTRICOS EN CARGA'' + tiempo máximo de permanencia\\
\end{tabular}};

% =============== D6: emergencia y mantenimiento
\node[tit] at (237,133) {D6 -- EMERGENCIA Y MANTENIMIENTO};
\node[anchor=north west,font=\fontsize{5.4}{6.6}\selectfont,align=left,text width=165mm] at (237,126) {%
\textbf{Ante incendio, choque o cable dañado:} 1) pulsar EPO-1 o EPO-2 (abre TD-VE); 2) llamar al 9-1-1 e informar ``estación de carga de vehículos eléctricos, 480 V''; 3) si el fuego es en la batería de un vehículo, no usar agua sobre el equipo energizado y esperar a Bomberos; 4) abrir IP-VE en el nicho del transformador si TD-VE no es accesible.\\[1mm]
\textbf{Extintor:} uno de CO$_2$ o polvo ABC de 10 lb junto a la celda eléctrica (RNPCI) (S).\\[1mm]
\textbf{Mantenimiento preventivo (S):} mensual: inspección visual de cables, conectores y fundas; prueba del botón de cada CCID/RCD. Semestral: limpieza de filtros de DC-1, apriete de terminales con termografía, prueba de EPO. Anual: medición de la resistencia de tierra, prueba de aislamiento y verificación de la firma del EMS. Contrato con el representante de Autel; repuestos críticos: un conector J1772 y un cable CCS1.\\[1mm]
\textbf{Ampliación:} T-VE admite 2 cargadores CA más (espacios de reserva en TD-AC); TD-VE tiene un espacio 3P de reserva y existe un ducto de 3'' de reserva para un segundo cargador CC, sujeto al estudio de capacidad del ICE.};

\cajetin{VE-03}{Detalles de montaje, zanja, base del cargador rápido,}{puesta a tierra, rotulación, emergencia y mantenimiento}{Indicadas}{4 de 4}
\end{tikzpicture}
\end{document}
